\chapter{Introduction}

Microorganisms are ubiquitous in the biosphere, contributing significantly to nutrient and carbon cycling. Biological processes in glaciers and ice sheets are dominated almost exclusively by microbial communities \cite{AnesioAm2017Tmog}. Arctic microbes contribute to the build-up of biomass and macronutrients such as carbon, phosphorus and nitrogen in soils \cite{BradleyJamesA2014Mcdi}, as well as impacting the fertilization and productivity of Polar fjords and oceans. They can also modify the albedo of ice surfaces, affecting the rate of sea-level rise. \\

\noindent To proliferate in extreme environments, microbes have developed survival strategies, including dormancy. Dormancy is a reversible state of low metabolic activity utilised by microorganisms when energetic resources become scarce \cite{BradleyJamesA.2019Sotf}. Dormant microorganisms generate a seed bank of individuals, ready for resuscitation when favourable environmental conditions return \cite{JayT.Lennon2011Msbt}. \\

\noindent Dormancy is significant in microbial biogeography because it reduces dispersal limitation (stochastic mechanism) and resource competition (deterministic mechanism) \cite{locey2010synthesizing}. Dormancy must then be considered when quantifying the contributions of deterministic and stochastic mechanisms to microbial biogeography and its implication for a changing polar habitat. Despite this, little is known about the interplay between biogeography and dormancy.\\

\noindent Understanding the microbial ecology and biogeography of cryoenvironments is essential as both poles experience rapid climate change, affecting both soil and glacial systems \cite{SicilianoStevenD2014Sfia}. Historical records show that Arctic air temperature has warmed twice as fast as the global mean and may warm by more than 4\degree C by the end of the 21st Century \cite{stocker2014climate}. The adaptation of microorgamisms to extreme polar environments makes them among the most vulnerable to climate change \cite{WarwickFVincent2010Mert}. \noindent Previous research has shown that the recent rapid retreat of glaciers as a result of global warming exposes Arctic soils to microbial colonisation \cite{10.3389/feart.2017.00026}. As temperatures rise in the Arctic, microbial decomposition of the active layer and permafrost soil carbon could positively affect warming. \\ 

\noindent This project will quantify the activity and dormancy of microbes in Arctic glaciers and soils in a biogeographical framework. I will conduct fieldwork to investigate the spatial determinants of microbial community patterns. I will incorporate novel methods to access the active state of those communities. Finally I will implement this empirical data in developing a predictive numerical model, integrating biogeography and dormancy.\\
